\newpage
\vspace{0.8cm}

\section{Ordine del giorno}

\subsection{Allineamento tecnico sul boilerplate del PoC}
Nel corso della riunione è stato presentato e discusso il codice del boilerplate predisposto per il Proof of Concept (PoC), illustrando il relativo \texttt{README} e lo stack adottato (Next.js con App Router, \texttt{pnpm}, TypeScript, Drizzle ORM e Supabase):
\begin{itemize}
    \item \textbf{Ambiente locale con Supabase e Docker}: il backend Supabase viene avviato in locale tramite Docker per lavorare in ambiente isolato senza dipendere dall'istanza online; è stata inoltre segnalata la mancanza del file \texttt{.env.example} in \texttt{apps/web}, da aggiungere al repository;
    \item \textbf{Autenticazione Google OAuth}: il login con Google è gestito tramite Supabase; le credenziali di sviluppo configurate su account personale non dovranno essere impiegate in produzione e la protezione delle rotte avviene tramite middleware applicativo;
    \item \textbf{Database e Row-Level Security (RLS)}: gli schemi e le migrazioni sono gestiti con Drizzle ORM, mentre per tutte le query applicative è richiesto l'utilizzo del wrapper dedicato (\texttt{withRLS}) al fine di garantire il rispetto delle policy di Row-Level Security;
    \item \textbf{Video walkthrough}: la registrazione della presentazione tecnica del boilerplate verrà caricata sul canale YouTube del gruppo (ad uso esclusivamente interno e accessibile solo tramite link diretto) a cura di F. Lissandrin e condivisa con i membri del team.
\end{itemize}

\subsection{Account Vercel Pro, pipeline CI/CD e gestione degli account GitHub}
Il gruppo ha ottenuto l'accesso all'account Vercel Pro messo a disposizione dall'azienda, riscontrando tuttavia alcune criticità tecniche nella configurazione della pipeline:
\begin{itemize}
    \item \textbf{Fallimento del deploy}: manca il file di configurazione (\texttt{vercel.config.js} oppure \texttt{vercel.config.ts}) all'interno del repository, necessario per mantenere la configurazione allineata alla codebase ad ogni push;
    \item \textbf{Problematica legata agli account GitHub}: i push effettuati verso il repository collegato a Vercel da account GitHub personali potrebbero comportare il conteggio di membri aggiuntivi a pagamento sul piano team. Come possibile soluzione tecnica alternativa è stato ipotizzato un repository mirror; il tema verrà comunque chiarito direttamente con il referente aziendale nell'incontro di domani o verificato tramite test mirato;
    \item \textbf{Completamento pipeline}: la sistemazione della pipeline Vercel e dell'integrazione con GitHub è stata affidata congiuntamente a F. Lissandrin e A. Tang, con obiettivo di completamento entro la fine dello Sprint 1 (24/11/2026).
\end{itemize}

\newpage
\subsection{Avanzamento della documentazione di progetto e automazione del Glossario}
È stato fatto il punto sullo stato della documentazione in vista della scadenza di fine Sprint 1 (24/11/2026):
\begin{itemize}
    \item \textbf{Norme di Progetto}: il documento dovrà essere ulteriormente integrato a cura di G. Gambarro con le nuove procedure e convenzioni definite dal gruppo, in particolare la struttura delle cartelle del repository e l'aggiunta di \emph{Granola} tra gli strumenti utilizzati per le trascrizioni delle riunioni;
    \item \textbf{Glossario e automazione della marcatura}: è stato ribadito che i termini presenti nel Glossario devono essere contrassegnati con la lettera $G$ a pedice su \textbf{ogni occorrenza} in tutti i documenti (non soltanto alla prima comparsa), introducendo inoltre l'iniziale alfabetica e il collegamento cliccabile. Per evitare errori e appesantimenti manuali, è stata progettata l'introduzione di uno script automatico (affidato ad A. Tang) che legga i termini del Glossario, li individui nei sorgenti e inserisca il markup LaTeX (\texttt{\textbackslash gls}) tra il repository documentale e la pubblicazione sul sito web, così da aggiornare anche i documenti già redatti;
    \item \textbf{Piano di Progetto}: G. Grinovero proseguirà la stesura delle sezioni 3 e 4 (modello di sviluppo e preventivo) sulla base delle informazioni attualmente disponibili, mentre il Responsabile (C. Antoniazzo) redigerà la previsione oraria per lo Sprint 1 e, a conclusione dello sprint, compilerà il consuntivo di periodo confrontando ore previste ed effettive;
    \item \textbf{Analisi dei Requisiti}: il gruppo ha concordato l'avvio del documento con la creazione della cartella dedicata su GitHub e l'apertura di un branch secondo la convenzione standard. D. Grosselle curerà l'impostazione della struttura, l'introduzione, la raccolta delle definizioni dalle slide del Prof. Cardin e l'avvio dei primi casi d'uso a seguito dei chiarimenti con il proponente.
\end{itemize}

\subsection{Preparazione dell'incontro esterno con Halue (18 novembre 2026)}
In vista dell'incontro esterno con il proponente Halue fissato per il 18/11/2026 alle ore 13:30, il team ha definito l'agenda operativa della demo e le domande chiave da sottoporre:
\begin{itemize}
    \item presentare il PoC relativo alla pagina di login;
    \item chiarire se per i push su GitHub e le Pull Request verso la pipeline Vercel sia necessario utilizzare un account condiviso oppure se possano essere impiegati gli account personali dei membri del team;
    \item definire con precisione quali visualizzazioni, pagine, statistiche, elementi da escludere e pulsanti con azioni specifiche siano attesi nella dashboard, così da derivare correttamente i casi d'uso.
\end{itemize}
La redazione del Verbale Esterno dell'incontro del 18/11/2026 è stata assegnata a G. Gambarro, con consegna prevista entro due giorni (scadenza 20/11/2026).

\section{Programmazione prossimo incontro}
L'incontro esterno con l'azienda proponente Halue è programmato per il giorno \textbf{18/11/2026 alle ore 13:30}, occasione in cui verrà mostrato il PoC di login e verranno discussi i chiarimenti sui casi d'uso della dashboard e sulla gestione GitHub/Vercel.

\section{Tabella decisioni prese}
\renewcommand{\arraystretch}{1.3}
\noindent
\begin{tabularx}{\textwidth}{|l|X|}
    \hline
    \textbf{Codice} & \textbf{Decisione} \\
    \hline
    VI\_2026-11-17\_D01 & Allineamento e adozione del boilerplate del PoC (Next.js, Drizzle ORM, Supabase locale via Docker) con utilizzo del wrapper \texttt{withRLS} e condivisione interna del video walkthrough. \\
    \hline
    VI\_2026-11-17\_D02 & Adozione dell'account Vercel Pro fornito da Halue e pianificazione dell'inserimento di \texttt{vercel.config.js/ts} e \texttt{.env.example} per completare la pipeline entro il 24/11/2026. \\
    \hline
    VI\_2026-11-17\_D03 & Estensione della marcatura del Glossario ($G$ a pedice e link cliccabile) a ogni occorrenza in tutti i documenti mediante sviluppo di uno script di automazione. \\
    \hline
    VI\_2026-11-17\_D04 & Aggiornamento delle Norme di Progetto con la documentazione della struttura del repository e l'introduzione di Granola tra gli strumenti di trascrizione adottati dal gruppo. \\
    \hline
    VI\_2026-11-17\_D05 & Avvio dell'Analisi dei Requisiti (struttura, introduzione, definizioni dalle slide del Prof. Cardin e primi casi d'uso) e prosecuzione del Piano di Progetto (sezioni 3, 4 e ore preventivate Sprint 1) con scadenza 24/11/2026. \\
    \hline
    VI\_2026-11-17\_D06 & Definizione della scaletta e delle domande da porre ad Halue nell'incontro del 18/11/2026 alle ore 13:30 (demo PoC login, account GitHub/Vercel e visualizzazioni della dashboard per i casi d'uso), con redazione del verbale esterno entro il 20/11/2026. \\
    \hline
\end{tabularx}
